% =============================================================================
%  Rapport de Projet de Fin d'Annee (PFA)
%  Detection intelligente des fraudes dans les systemes financiers et bancaires
%  basee sur l'integration de l'IA et des Systemes Multi-Agents
% =============================================================================
\documentclass[12pt,a4paper,oneside]{report}

% ---------------------------- PACKAGES ---------------------------------------
\usepackage[utf8]{inputenc}
\usepackage[T1]{fontenc}
\usepackage[french]{babel}
\usepackage{lmodern}
\usepackage[a4paper,margin=2.5cm]{geometry}
\usepackage{graphicx}
\usepackage{xcolor}
\usepackage{hyperref}
\usepackage{amsmath,amssymb}
\usepackage{booktabs}
\usepackage{array}
\usepackage{tabularx}
\usepackage{longtable}
\usepackage{enumitem}
\usepackage{fancyhdr}
\usepackage{titlesec}
\usepackage{listings}
\usepackage{caption}
\usepackage{subcaption}
\usepackage{tikz}
\usetikzlibrary{arrows.meta,positioning,shapes.geometric,shadows,matrix,fit}
\usepackage{pgfplots}
\pgfplotsset{compat=1.18}
\usepgfplotslibrary{statistics}
\usepackage{float}
\usepackage[backend=biber,style=numeric,sorting=none]{biblatex}

% ---------------------------- COULEURS ---------------------------------------
\definecolor{pfaBlue}{RGB}{30,90,160}
\definecolor{pfaDark}{RGB}{20,40,80}
\definecolor{pfaGray}{RGB}{120,120,120}
\definecolor{codeBg}{RGB}{248,248,248}
\definecolor{codeKey}{RGB}{0,90,180}
\definecolor{codeStr}{RGB}{160,40,40}
\definecolor{codeCmt}{RGB}{80,140,80}

% ---------------------------- HYPERREF ---------------------------------------
\hypersetup{
  colorlinks=true,
  linkcolor=pfaDark,
  urlcolor=pfaBlue,
  citecolor=pfaBlue,
  pdftitle={Detection intelligente de fraudes - PFA},
  pdfauthor={Equipe PFA}
}

% ---------------------------- STYLES TITRES ----------------------------------
\titleformat{\chapter}[display]
  {\normalfont\Huge\bfseries\color{pfaDark}}
  {\filleft\Large Chapitre~\thechapter}
  {1ex}
  {\titlerule\vspace{1ex}\filright}
  [\vspace{1ex}\titlerule]

\titleformat{\section}{\Large\bfseries\color{pfaDark}}{\thesection}{1em}{}
\titleformat{\subsection}{\large\bfseries\color{pfaBlue}}{\thesubsection}{1em}{}

% ---------------------------- ENTETES ----------------------------------------
\pagestyle{fancy}
\fancyhf{}
\fancyhead[L]{\small\itshape\leftmark}
\fancyhead[R]{\small PFA $\bullet$ Detection de fraude}
\fancyfoot[C]{\thepage}
\renewcommand{\headrulewidth}{0.4pt}

% ---------------------------- LISTINGS ---------------------------------------
\lstdefinestyle{codeStyle}{
  backgroundcolor=\color{codeBg},
  basicstyle=\ttfamily\small,
  keywordstyle=\color{codeKey}\bfseries,
  stringstyle=\color{codeStr},
  commentstyle=\color{codeCmt}\itshape,
  numbers=left,
  numberstyle=\tiny\color{pfaGray},
  numbersep=8pt,
  frame=single,
  rulecolor=\color{pfaGray!40},
  breaklines=true,
  breakatwhitespace=true,
  showstringspaces=false,
  tabsize=2,
  captionpos=b,
  xleftmargin=12pt,
  framexleftmargin=10pt,
}
\lstset{style=codeStyle}

% ---------------------------- BIBLIOGRAPHIE ----------------------------------
\addbibresource{references.bib}

% ============================================================================
\begin{document}
% ============================================================================

% ----------------------------- PAGE DE GARDE --------------------------------
\begin{titlepage}
  \centering
  \vspace*{1cm}

  {\large\bfseries Universite [Nom de l'Universite]\par}
  {\large Faculte des Sciences \\ Departement d'Informatique\par}
  \vspace{0.6cm}

  \rule{\linewidth}{0.6pt}\\[0.4cm]
  {\Huge\bfseries\color{pfaDark} Detection intelligente des fraudes\\[0.2cm]
   dans les systemes financiers et bancaires\par}
  \vspace{0.4cm}
  {\Large\itshape basee sur l'integration de l'Intelligence Artificielle\\
   et des Systemes Multi-Agents\par}
  \rule{\linewidth}{0.6pt}\\[1cm]

  {\large\bfseries Projet de Fin d'Annee (PFA)\par}
  {\large Annee universitaire 2025--2026\par}
  \vspace{1.5cm}

  \begin{flushleft}
  {\large\bfseries Realise par :} \\[0.2cm]
  \begin{tabular}{ll}
    Etudiant 1 & --- Backend SMA JADE \\
    Etudiant 2 & --- Pipeline Machine Learning \\
    Etudiant 3 & --- Frontend \& UX \\
    Etudiant 4 & --- Donnees \& DevOps \\
  \end{tabular}
  \end{flushleft}
  \vspace{1cm}
  \begin{flushleft}
  {\large\bfseries Encadre par :} M. / Mme [Nom de l'encadrant]\par
  \end{flushleft}

  \vfill
  {\small\itshape Soutenu publiquement le \rule{3cm}{0.3pt}\par}
\end{titlepage}

% ----------------------------- DEDICACES ------------------------------------
\chapter*{Dedicaces}
\addcontentsline{toc}{chapter}{Dedicaces}

\noindent A nos chers parents, qui n'ont jamais cesse de nous soutenir tout au
long de notre cursus universitaire, et qui ont fait de nous ce que nous sommes
aujourd'hui. Aucune dedicace ne saurait exprimer notre reconnaissance, notre
amour, et le profond respect que nous leur portons.

\bigskip
\noindent A nos freres, soeurs, et tous les membres de nos familles, pour leur
soutien moral et leurs encouragements constants.

\bigskip
\noindent A nos amis et camarades de promotion, pour les moments partages, les
debats enrichissants et l'entraide qui ont marque ces annees d'etudes.

\bigskip
\noindent A tous nos enseignants, qui ont contribue a notre formation et nous
ont transmis le gout du savoir et de la rigueur.

\bigskip
\hfill {\itshape L'equipe PFA}

% ----------------------------- REMERCIEMENTS ---------------------------------
\chapter*{Remerciements}
\addcontentsline{toc}{chapter}{Remerciements}

Au terme de ce travail, il nous est agreable de nous acquitter d'une dette de
reconnaissance auprès de toutes les personnes dont l'intervention au cours de
ce projet a favorise son aboutissement.

Nous tenons en premier lieu a exprimer notre sincere gratitude a notre
encadrant, \textbf{[Nom de l'encadrant]}, pour la qualite de son encadrement,
sa disponibilite, sa patience, ses conseils avises et ses critiques
constructives qui nous ont permis de mener a bien ce projet.

Nous remercions egalement l'ensemble du corps professoral du
\textbf{Departement d'Informatique} de la \textbf{Faculte [...]} pour la qualite
de l'enseignement dispense durant notre cursus, ainsi que pour leur
disponibilite et leur ecoute.

Nos remerciements s'adressent enfin aux \textbf{membres du jury} qui nous ont
fait l'honneur d'evaluer ce travail, et a toute personne ayant contribue, de
près ou de loin, a la realisation de ce projet.

% ----------------------------- RESUME ---------------------------------------
\chapter*{Resume}
\addcontentsline{toc}{chapter}{Resume}

La fraude bancaire represente aujourd'hui un defi majeur pour les institutions
financieres, avec des pertes mondiales estimees a plusieurs dizaines de
milliards de dollars par an. Les approches classiques fondees uniquement sur
des regles statiques ou sur un modele d'apprentissage isole montrent leurs
limites face a la sophistication croissante des fraudeurs et au volume des
transactions a traiter en temps reel.

Ce projet de fin d'annee propose une architecture hybride combinant un
\textbf{systeme multi-agents} (SMA) developpe avec la plateforme \textbf{JADE}
et un \textbf{pipeline d'intelligence artificielle} reposant sur un
\emph{ensemble} \textbf{XGBoost} (apprentissage supervise) et
\textbf{Isolation Forest} (apprentissage non supervise). Six agents autonomes
collaborent via le standard \textbf{FIPA-ACL} et le protocole
\emph{FIPA Contract-Net} pour analyser, scorer, evaluer et trancher chaque
transaction. Trois regles metier (montant, velocite, horaire) viennent enrichir
le score machine learning afin de produire une decision \emph{explicable}.

Le systeme est integralement conteneurise (\textbf{Docker Compose}) et expose
ses alertes en temps reel a un \textbf{tableau de bord React} via
\textbf{WebSocket}. Toutes les decisions sont persistees dans une base
\textbf{PostgreSQL}. Les modeles ont ete entraines sur le jeu de donnees
public de Kaggle \emph{Credit Card Fraud Detection} ($284\,807$ transactions,
$492$ fraudes) et atteignent un \textbf{ROC-AUC d'ensemble de 0,96}.

\noindent\textbf{Mots-cles :} Detection de fraude, Systemes Multi-Agents, JADE,
FIPA-ACL, Machine Learning, XGBoost, Isolation Forest, Approche hybride,
Temps reel, WebSocket.

% ----------------------------- ABSTRACT --------------------------------------
\chapter*{Abstract}
\addcontentsline{toc}{chapter}{Abstract}

Banking fraud has become one of the most pressing challenges faced by financial
institutions, with global losses estimated at tens of billions of dollars
annually. Traditional approaches relying solely on static rules or on a single
isolated learning model are increasingly inadequate against modern adversaries
and the sheer volume of real-time transactions.

This Final Year Project proposes a hybrid architecture combining a
\textbf{Multi-Agent System} (MAS) developed using the \textbf{JADE} platform
with a \textbf{Machine Learning} pipeline based on an \emph{ensemble} of
\textbf{XGBoost} (supervised learning) and \textbf{Isolation Forest}
(unsupervised anomaly detection). Six autonomous agents cooperate through the
\textbf{FIPA-ACL} standard and the \emph{FIPA Contract-Net} protocol to
analyse, score, evaluate and decide on each transaction. Three business rules
(amount, velocity, timing) enrich the ML score to produce an
\emph{explainable} decision.

The system is fully containerised (\textbf{Docker Compose}) and broadcasts
real-time alerts to a \textbf{React dashboard} through \textbf{WebSocket}, with
full persistence in a \textbf{PostgreSQL} database. Models were trained on the
public Kaggle \emph{Credit Card Fraud Detection} dataset ($284{,}807$
transactions, $492$ frauds) and achieve a \textbf{ROC-AUC ensemble score of
0.96}.

\noindent\textbf{Keywords:} Fraud detection, Multi-Agent Systems, JADE,
FIPA-ACL, Machine Learning, XGBoost, Isolation Forest, Hybrid approach,
Real-time, WebSocket.

% ----------------------------- TABLE DES MATIERES ----------------------------
\tableofcontents
\listoffigures
\listoftables

% ============================================================================
%                          INTRODUCTION GENERALE
% ============================================================================
\chapter*{Introduction generale}
\addcontentsline{toc}{chapter}{Introduction generale}
\markboth{Introduction generale}{}

\section*{Contexte}

La numerisation des services bancaires et l'explosion du commerce electronique
ont profondement transforme les habitudes de consommation. En contrepartie,
elles ont ouvert un boulevard a une nouvelle generation de criminalite
financiere : la fraude bancaire. Selon le rapport \emph{Nilson Report} (2023),
les pertes mondiales liees a la fraude par carte ont depasse les
\textbf{32 milliards de dollars} en 2022, avec une projection a plus de
\textbf{40 milliards en 2027}. Plus encore, ces fraudes sont desormais
massivement automatisees, perpetrees a l'echelle internationale, et exploitent
des techniques toujours plus elaborees : \emph{card testing}, \emph{card not
present}, vol d'identite, ingenierie sociale\dots

Face a ce phenomene, les institutions financieres sont contraintes
d'industrialiser leur lutte contre la fraude. Or, les systemes legacy --
souvent fondes sur des batteries de regles statiques ecrites a la main --
peinent a suivre. D'une part, ils generent un taux excessif de
\emph{faux positifs} qui irrite les clients legitimes ; d'autre part, ils sont
incapables de detecter les nouvelles formes de fraude qui n'ont pas encore ete
caracterisees explicitement.

\section*{Problematique}

L'enjeu de notre projet est de repondre a la question suivante :

\begin{quote}
\itshape Comment concevoir un systeme de detection de fraude bancaire qui soit
a la fois \textbf{intelligent} (capable d'apprendre et de detecter des
patterns nouveaux), \textbf{explicable} (auditable par un humain et conforme
aux exigences reglementaires), \textbf{distribue} (capable de monter en charge
sur de gros volumes), et \textbf{temps reel} (decision en quelques
millisecondes) ?
\end{quote}

\section*{Objectifs}

Ce travail vise a concevoir et a implementer un prototype operationnel
combinant deux paradigmes complementaires :

\begin{itemize}[leftmargin=*]
  \item \textbf{L'Intelligence Artificielle (IA)} pour apprendre les patterns
        de fraude a partir de donnees historiques (apprentissage supervise) et
        detecter des anomalies inedites (apprentissage non supervise).
  \item \textbf{Les Systemes Multi-Agents (SMA)} pour orchestrer la chaine de
        decision de maniere distribuee, modulaire, et conforme aux standards
        FIPA-ACL.
\end{itemize}

Plus precisement, le systeme doit :
\begin{enumerate}[leftmargin=*]
  \item Recevoir un flux continu de transactions bancaires en temps reel ;
  \item Calculer pour chaque transaction un score de risque combinant un
        modele ML d'ensemble et un ensemble de regles metier ;
  \item Trancher chaque transaction en \texttt{APPROUVER}, \texttt{SIGNALER}
        ou \texttt{BLOQUER} ;
  \item Persister chaque decision et chaque message ACL pour audit ;
  \item Diffuser les alertes vers un tableau de bord operationnel en temps
        reel.
\end{enumerate}

\section*{Organisation du rapport}

Ce rapport est structure en six chapitres :

\begin{itemize}[leftmargin=*]
  \item Le \textbf{chapitre 1} presente le contexte general, la problematique
        de la fraude bancaire et un etat de l'art des techniques de detection.
  \item Le \textbf{chapitre 2} detaille les concepts theoriques mobilises :
        Systemes Multi-Agents, FIPA-ACL, plateforme JADE, et algorithmes
        d'apprentissage automatique utilises (XGBoost, Isolation Forest, SMOTE).
  \item Le \textbf{chapitre 3} presente les frameworks, langages de
        programmation et outils selectionnes pour la realisation.
  \item Le \textbf{chapitre 4} expose la conception du systeme : architecture
        globale, modele d'agents, modele de donnees, et protocole d'interaction.
  \item Le \textbf{chapitre 5} decrit l'implementation technique : organisation
        du code, agents JADE, microservice ML et choix d'ingenierie.
  \item Le \textbf{chapitre 6} presente les tests, les resultats experimentaux
        et les captures d'ecran du tableau de bord.
\end{itemize}

Une \textbf{conclusion generale} dresse le bilan du projet et trace plusieurs
pistes d'amelioration et de recherche future.

% ============================================================================
%                          CHAPITRE 1 - CONTEXTE
% ============================================================================
\chapter{Contexte general et etat de l'art}

\section{Introduction du chapitre}
Ce chapitre pose les fondations conceptuelles et factuelles du projet. Apres
avoir cerne le phenomene de la fraude bancaire, nous passons en revue les
principales familles d'approches utilisees aujourd'hui pour la detecter, en
soulignant leurs forces et leurs limites. Cette analyse motive le choix d'une
architecture hybride combinant IA et SMA.

\section{La fraude bancaire : un phenomene en pleine mutation}

\subsection{Definition et typologie}

La \emph{fraude bancaire} designe toute action illegale visant a obtenir, par
tromperie, un gain financier au detriment d'un etablissement bancaire ou d'un
client. On distingue communement :

\begin{itemize}[leftmargin=*]
  \item \textbf{La fraude par carte} (CNP, \emph{Card Not Present}, vol
        physique, skimming, clonage) ;
  \item \textbf{La fraude par virement} (manipulation du beneficiaire,
        usurpation d'identite, faux ordre du president) ;
  \item \textbf{La fraude au paiement instantane} (\emph{authorized push
        payment}) ;
  \item \textbf{La cyberfraude} (phishing, malware bancaire, attaques par
        injection sur l'API).
\end{itemize}

\subsection{Caracteristiques statistiques}

Les fraudes presentent quelques caracteristiques recurrentes qui guident les
strategies de detection :

\begin{itemize}[leftmargin=*]
  \item \textbf{Tres faible prevalence} : moins de 0,2\,\% des transactions
        sont frauduleuses dans les jeux de donnees publics.
  \item \textbf{Patterns temporels} : les fraudes par rafale (\emph{burst
        attack}) tentent de maximiser le gain en quelques minutes avant
        detection.
  \item \textbf{Horaires atypiques} : une part importante survient hors
        heures ouvrees ou la nuit (UTC).
  \item \textbf{Adaptabilite} : les fraudeurs modifient leurs patterns des
        qu'une regle est detectee, rendant les systemes purement
        deterministes obsoletes.
\end{itemize}

\section{Etat de l'art des techniques de detection}

\subsection{Approches a base de regles (rule-based)}

Historiquement, les premiers systemes anti-fraude se sont appuyes sur des
batteries de regles ecrites par des experts metier : \emph{si le montant
depasse X et que le pays est dans la liste Y alors bloquer}. Ces approches
sont :

\begin{itemize}[leftmargin=*]
  \item[\(+\)] simples, transparentes et auditables ;
  \item[\(+\)] rapides a executer et faciles a deployer ;
  \item[\(-\)] rigides, incapables de detecter de nouvelles formes de fraude ;
  \item[\(-\)] generatrices d'un taux eleve de faux positifs ;
  \item[\(-\)] couteuses a maintenir a mesure que les regles s'accumulent.
\end{itemize}

\subsection{Approches statistiques classiques}

L'analyse statistique (regression logistique, scorecards, analyse
discriminante) a longtemps constitue le standard en \emph{credit risk}. Elle
ameliore la rigueur mathematique des regles mais reste limitee face a des
donnees a forte dimension et a relations non lineaires.

\subsection{Apprentissage automatique supervise}

Avec la disponibilite croissante de donnees etiquetees, les algorithmes
supervises se sont imposes :

\begin{itemize}[leftmargin=*]
  \item \textbf{Random Forest} : robuste, peu sensible au parametrage,
        explicabilite via importance des features ;
  \item \textbf{Gradient Boosting (XGBoost, LightGBM)} : etat de l'art en
        precision sur donnees tabulaires, gere nativement le desequilibre via
        le parametre \texttt{scale\_pos\_weight} ;
  \item \textbf{Reseaux de neurones} : interessants pour les sequences
        (LSTM/Transformers) mais moins performants sur des features tabulaires
        statiques.
\end{itemize}

\subsection{Apprentissage non supervise et detection d'anomalies}

Pour detecter les fraudes \emph{nouvelles} (jamais vues dans les donnees
d'entrainement), les methodes non supervisees jouent un role complementaire :

\begin{itemize}[leftmargin=*]
  \item \textbf{Isolation Forest} : isole les points anormaux au moyen d'arbres
        aleatoires.
  \item \textbf{One-Class SVM} : modelise la frontiere des points "normaux".
  \item \textbf{Autoencoders} : exploitent la reconstruction pour detecter
        l'anomalie.
\end{itemize}

\subsection{Approches hybrides}

La litterature recente \cite{west2016intelligent,abakarim2018hybrid} converge
vers des approches \emph{hybrides} combinant :

\begin{itemize}[leftmargin=*]
  \item Un modele supervise (XGBoost) pour les fraudes connues ;
  \item Un detecteur d'anomalies (Isolation Forest) pour les fraudes nouvelles ;
  \item Des regles metier deterministes pour la conformite reglementaire et
        l'explicabilite ;
  \item Une orchestration multi-agents pour la modularite et la
        distribution.
\end{itemize}

C'est cette derniere approche hybride que nous avons choisi d'implementer.

\section{Synthese et positionnement}

\begin{table}[H]
\centering
\caption{Comparaison des principales approches de detection de fraude}
\label{tab:comparaison-approches}
\renewcommand{\arraystretch}{1.3}
\begin{tabularx}{\textwidth}{l X X X}
\toprule
\textbf{Approche} & \textbf{Avantages} & \textbf{Limites} & \textbf{Cas d'usage} \\
\midrule
Regles & Simple, explicable & Rigide, faux positifs & Conformite \\
Statistique & Rigueur math. & Faible avec features non lineaires & Scoring credit \\
ML supervise & Tres precis & Necessite labels & Fraudes connues \\
ML non supervise & Detecte nouveau & Faux positifs eleves & Anomalies \\
SMA & Modulaire, distribue & Complexite & Orchestration \\
\textbf{Hybride (notre approche)} & \textbf{Combine tous les avantages} & \textbf{Plus complexe} & \textbf{Production} \\
\bottomrule
\end{tabularx}
\end{table}

\section{Conclusion du chapitre}
Ce chapitre a montre que la lutte contre la fraude bancaire ne peut plus se
contenter d'une approche unique. La voie hybride combinant IA et SMA,
plebiscitee par la litterature recente, est celle que nous explorons dans la
suite de ce rapport.

% ============================================================================
%                          CHAPITRE 2 - CONCEPTS
% ============================================================================
\chapter{Concepts theoriques}

\section{Introduction du chapitre}
Ce chapitre presente les concepts theoriques mobilises dans notre solution :
les Systemes Multi-Agents (SMA), le standard FIPA-ACL, la plateforme JADE,
ainsi que les algorithmes d'apprentissage automatique XGBoost, Isolation
Forest et la technique de reequilibrage SMOTE.

\section{Les Systemes Multi-Agents (SMA)}

\subsection{Notion d'agent}

Selon la definition de Wooldridge \cite{wooldridge2009introduction}, un
\textbf{agent} est une entite logicielle situee dans un environnement, capable
d'agir de maniere autonome pour atteindre ses objectifs. Un agent presente
typiquement les proprietes suivantes :

\begin{itemize}[leftmargin=*]
  \item \textbf{Autonomie} : il agit sans intervention humaine directe ;
  \item \textbf{Reactivite} : il percoit son environnement et reagit aux
        changements ;
  \item \textbf{Pro-activite} : il poursuit ses propres buts ;
  \item \textbf{Sociabilite} : il interagit avec d'autres agents.
\end{itemize}

\subsection{Systemes Multi-Agents}

Un \textbf{Systeme Multi-Agents} (SMA) est un ensemble d'agents qui cohabitent
dans un environnement et cooperent (ou se concurrencent) pour atteindre des
buts individuels et collectifs. Les SMA sont particulierement adaptes a des
problemes :

\begin{itemize}[leftmargin=*]
  \item \emph{Distribues} : la connaissance et le controle sont repartis ;
  \item \emph{Dynamiques} : l'environnement evolue en permanence ;
  \item \emph{Modulaires} : le systeme se decompose en services
        independants.
\end{itemize}

La detection de fraude bancaire en temps reel coche les trois cases.

\subsection{Standard FIPA-ACL}

La \emph{Foundation for Intelligent Physical Agents} (FIPA) a normalise un
langage de communication entre agents : le \textbf{ACL} (\emph{Agent
Communication Language}). Un message ACL contient notamment :

\begin{itemize}[leftmargin=*]
  \item un \textbf{performatif} (intention) tel que \texttt{REQUEST},
        \texttt{INFORM}, \texttt{PROPOSE}, \texttt{ACCEPT\_PROPOSAL},
        \texttt{REJECT\_PROPOSAL}, \texttt{REFUSE}, \texttt{FAILURE} ;
  \item un \textbf{emetteur} et un ou plusieurs \textbf{destinataires} ;
  \item une \textbf{ontologie} qui definit le vocabulaire metier ;
  \item un \textbf{contenu} (typiquement du JSON ou du SL) ;
  \item un \textbf{identifiant de conversation} permettant de correler
        plusieurs messages.
\end{itemize}

\subsection{Le protocole FIPA Contract-Net}

Le protocole \emph{Contract-Net} \cite{smith1980contract} est un standard
FIPA permettant la repartition dynamique de taches. Il repose sur quatre
phases :

\begin{enumerate}[leftmargin=*]
  \item Un agent \emph{initiator} envoie un \texttt{CFP} (Call For Proposal) ;
  \item Les \emph{participants} repondent par un \texttt{PROPOSE} ;
  \item L'\emph{initiator} accepte ou rejette via \texttt{ACCEPT\_PROPOSAL}
        ou \texttt{REJECT\_PROPOSAL} ;
  \item Le participant retenu execute la tache et retourne le resultat
        (\texttt{INFORM}).
\end{enumerate}

Dans notre projet, nous utilisons une version legerement adaptee : l'agent
Risque envoie directement un \texttt{PROPOSE} a l'agent Banque qui repond par
\texttt{ACCEPT\_PROPOSAL} (fraude confirmee) ou \texttt{REJECT\_PROPOSAL}
(transaction approuvee).

\subsection{La plateforme JADE}

\textbf{JADE} (\emph{Java Agent DEvelopment framework}) est une plateforme
open-source developpee par Telecom Italia Lab. Elle implemente fidelement les
specifications FIPA et fournit :

\begin{itemize}[leftmargin=*]
  \item Un \textbf{container} (conteneur d'agents) gerant le cycle de vie ;
  \item Un \textbf{Directory Facilitator (DF)} (sorte d'annuaire jaune) pour
        la decouverte dynamique d'agents par type de service ;
  \item Un \textbf{Agent Management System (AMS)} pour la gestion globale ;
  \item Une API riche de \emph{behaviours} (cyclic, one-shot, ticker, parallel,
        sequential\dots).
\end{itemize}

\section{Algorithmes d'apprentissage automatique}

\subsection{XGBoost}

\textbf{XGBoost} (\emph{eXtreme Gradient Boosting}) \cite{chen2016xgboost} est
une implementation hautement optimisee de l'algorithme de \emph{Gradient
Boosting Decision Trees}. L'idee centrale est d'entrainer une succession
d'arbres de decision faibles ou\`u chaque nouvel arbre corrige les erreurs des
precedents. Mathematiquement, on minimise une fonction objective :

\begin{equation}
  \mathcal{L}(\phi) = \sum_{i} \ell(\hat{y}_i, y_i) + \sum_{k} \Omega(f_k)
\end{equation}

ou $\ell$ est une fonction de perte (log-loss en classification binaire) et
$\Omega(f_k) = \gamma T + \frac{1}{2}\lambda \lVert w \rVert^2$ un terme de
regularisation sur la complexite de l'arbre $k$.

Pour gerer le desequilibre extreme entre classes, XGBoost permet de definir le
parametre \texttt{scale\_pos\_weight} qui pondere la classe minoritaire.

\subsection{Isolation Forest}

\textbf{Isolation Forest} \cite{liu2008isolation} est un algorithme non
supervise de detection d'anomalies. Son principe est simple : un point
\emph{anormal} est plus facilement \emph{isole} (separe des autres) qu'un
point \emph{normal}. L'algorithme construit une foret d'arbres aleatoires ou
chaque noeud divise l'espace selon un attribut et un seuil tires au hasard.

Le \emph{score d'anomalie} est inversement proportionnel a la profondeur
moyenne a laquelle un point est isole. On le normalise typiquement dans
$[0, 1]$ ou $1$ signifie tres anormal.

\subsection{SMOTE}

Sur le jeu de donnees Kaggle utilise, seulement $0{,}17\,\%$ des transactions
sont frauduleuses. Un classifieur naif obtiendrait $99{,}83\,\%$
d'\emph{accuracy} en predisant \og non-fraude\fg{} pour tout le monde.

Pour pallier ce probleme, nous utilisons \textbf{SMOTE} (\emph{Synthetic
Minority Over-sampling Technique}) \cite{chawla2002smote}. Cet algorithme
genere des exemples synthetiques de la classe minoritaire en interpolant
lineairement entre des voisins proches. Nous reequilibrons a un ratio
$1{:}10$ (10\,\% de fraudes apres SMOTE) afin de preserver le realisme tout
en aidant l'apprentissage.

\subsection{Ensemble XGBoost + Isolation Forest}

Les deux modeles sont complementaires :

\begin{itemize}[leftmargin=*]
  \item \emph{XGBoost} excelle sur les fraudes \og connues \fg{} (deja
        etiquetees dans le passe) ;
  \item \emph{Isolation Forest} repere les fraudes \og nouvelles \fg{} qui
        n'ont jamais ete vues mais s'ecartent statistiquement du normal.
\end{itemize}

Nous fusionnons les deux scores par une combinaison lineaire :

\begin{equation}
  \text{score}_{\text{ML}} = \alpha \cdot \text{score}_{\text{XGB}} +
                              (1 - \alpha) \cdot \text{score}_{\text{IF}}
  \quad\text{avec}~\alpha = 0{,}65.
\end{equation}

\section{Conclusion du chapitre}
Les concepts presentes ici fournissent les briques theoriques de notre
solution. Le chapitre suivant decrit comment ces briques sont assemblees dans
une architecture coherente.

% ============================================================================
%                CHAPITRE 3 - FRAMEWORKS, LANGAGES ET OUTILS
% ============================================================================
\chapter{Frameworks, langages et outils utilises}

\section{Introduction du chapitre}
Avant de detailler la conception, ce chapitre presente l'ensemble des
langages de programmation, frameworks et outils mobilises dans le projet. Ces
choix techniques ont ete dictes par trois exigences : la \textbf{maturite} du
framework (production-ready), sa \textbf{compatibilite} avec les standards
industriels, et la \textbf{communaute} active qui en assure la perennite.

\section{Langages de programmation}

\subsection{Java 17}

\textbf{Java} est un langage oriente objet, fortement type, compile en
\emph{bytecode} et execute par la JVM (\emph{Java Virtual Machine}). Il est
particulierement adapte aux systemes distribues critiques grace a sa
robustesse, son ecosysteme mature, et sa portabilite (\emph{write once, run
anywhere}). La version 17 (LTS) apporte des fonctionnalites modernes :
\emph{records}, \emph{sealed classes}, \emph{pattern matching},
\emph{switch expressions}\dots

\textbf{Role dans le projet :} implementation des six agents JADE, du serveur
WebSocket Tyrus, et de la couche de persistance Hibernate.

\subsection{Python 3.11}

\textbf{Python} est un langage interprete, dynamiquement type, plebiscite
dans la communaute Data Science / Machine Learning grace a son ecosysteme
extraordinairement riche (NumPy, Pandas, scikit-learn, TensorFlow,
PyTorch\dots). Sa syntaxe lisible accelere considerablement les iterations
de prototypage.

\textbf{Role dans le projet :} pipeline de feature engineering, entrainement
des modeles XGBoost et Isolation Forest, microservice Flask exposant
l'endpoint \texttt{/predict}.

\subsection{TypeScript 5}

\textbf{TypeScript} est un sur-ensemble de JavaScript ajoutant un systeme de
typage statique optionnel. Developpe par Microsoft, il s'est impose comme le
standard de fait pour les applications front-end de taille moyenne ou
grande, en eliminant a la compilation une grande partie des bugs typiques de
JavaScript.

\textbf{Role dans le projet :} developpement du tableau de bord React, avec
typage strict des messages WebSocket et des transitions d'etat Zustand.

\subsection{SQL (PostgreSQL)}

Le langage \textbf{SQL} (\emph{Structured Query Language}) reste le standard
universel pour l'interrogation des bases de donnees relationnelles. Notre
implementation utilise les extensions specifiques de PostgreSQL : type
\texttt{UUID} natif, contraintes \texttt{CHECK}, index sur \texttt{JSONB},
vues materialisees pour les KPI.

\textbf{Role dans le projet :} script de creation \texttt{init.sql},
requ\^etes Hibernate, vues d'agregation (\texttt{v\_kpi\_fraude}).

\section{Frameworks et bibliotheques}

\subsection{JADE (Java Agent DEvelopment framework)}

\textbf{JADE} \cite{jade2007programmer} est la plateforme multi-agents de
reference. Developpee par Telecom Italia Lab depuis 2000 et maintenue sous
licence LGPL, elle implemente fidelement les specifications FIPA :

\begin{itemize}[leftmargin=*]
  \item \textbf{Containers} : conteneurs d'agents avec gestion du cycle de
        vie (creation, demarrage, suspension, arr\^et) ;
  \item \textbf{Behaviours} : abstraction du comportement d'un agent
        (CyclicBehaviour, OneShotBehaviour, TickerBehaviour\dots) ;
  \item \textbf{Directory Facilitator (DF)} : annuaire \emph{pages jaunes}
        permettant aux agents de s'enregistrer et se decouvrir par type de
        service ;
  \item \textbf{Agent Management System (AMS)} : registre des agents et de
        leur statut ;
  \item \textbf{Message Transport Protocol} : transport des messages ACL via
        JICP, avec serialisation Java native.
\end{itemize}

\subsection{Apache Maven 3.9}

\textbf{Maven} est le gestionnaire de dependances et l'outil de build de
reference pour les projets Java. Il permet de declarer les dependances dans
un fichier \texttt{pom.xml} et automatise la compilation, l'execution des
tests, et la production d'artefacts (JAR, WAR\dots).

\textbf{Role dans le projet :} build du backend JADE et du simulateur, avec
le plugin \emph{maven-shade-plugin} pour produire des JAR \emph{fat}
(toutes dependances incluses).

\subsection{Hibernate 6 et JPA}

\textbf{Hibernate} est l'implementation de reference de la specification JPA
(\emph{Java Persistence API}). Il assure le mapping objet-relationnel
(\emph{ORM}) : les entites Java annotees \texttt{@Entity} sont
automatiquement persistees dans des tables SQL.

\textbf{Role dans le projet :} mapping des entites \texttt{Transaction},
\texttt{Alerte}, \texttt{DecisionAgent}, \texttt{MessageACL}.

\subsection{HikariCP}

\textbf{HikariCP} est le pool de connexions JDBC le plus performant
disponible (mesure 10\,000$\times$ plus rapide que les alternatives sur
certains benchmarks). Il gere efficacement la reutilisation des connexions
PostgreSQL.

\subsection{Tyrus + Grizzly}

\textbf{Tyrus} est l'implementation de reference de la specification
Jakarta WebSocket. Combine au serveur HTTP \textbf{Grizzly}, il fournit un
serveur WebSocket complet en quelques lignes de Java :

\begin{lstlisting}[language=Java,caption={Demarrage du serveur Tyrus},label=lst:tyrus]
Server tyrusServer = new Server("0.0.0.0", 8080, "/", null,
                                AlerteWebSocketServer.class);
tyrusServer.start();
\end{lstlisting}

\subsection{Flask 3}

\textbf{Flask} est un micro-framework web Python ultra-leger. Il fournit
juste ce qu'il faut pour exposer un endpoint REST sans imposer de structure
contraignante. Compare a Django, c'est l'outil ideal pour un microservice ML.

\textbf{Role dans le projet :} expose les endpoints \texttt{/predict},
\texttt{/health}, \texttt{/metrics} de \texttt{ml-api}.

\subsection{XGBoost 2.0}

\textbf{XGBoost} (\emph{eXtreme Gradient Boosting}) \cite{chen2016xgboost}
est une bibliotheque optimisee pour l'apprentissage par boosting d'arbres
de decision. Sa popularite est immense en competition Kaggle ou elle a
remporte des centaines de challenges.

\textbf{Caracteristiques :}
\begin{itemize}[leftmargin=*]
  \item Implementation native en C++ avec wrappers Python/R/JVM ;
  \item Support GPU \texttt{tree\_method=gpu\_hist} ;
  \item Gestion native des valeurs manquantes ;
  \item Regularisation L1/L2 integree ;
  \item Parametre \texttt{scale\_pos\_weight} pour le desequilibre.
\end{itemize}

\subsection{scikit-learn 1.5}

\textbf{scikit-learn} est la bibliotheque de Machine Learning de reference
en Python. Elle fournit une API unifiee \texttt{fit/predict} pour des
dizaines d'algorithmes (SVM, Random Forest, K-Means, PCA\dots).

\textbf{Role dans le projet :} \texttt{IsolationForest},
\texttt{train\_test\_split}, \texttt{classification\_report},
\texttt{roc\_auc\_score}, \texttt{confusion\_matrix}.

\subsection{imbalanced-learn (SMOTE)}

\textbf{imbalanced-learn} etend scikit-learn avec des methodes specifiques
au desequilibre de classes : SMOTE, ADASYN, sous-echantillonnage,
combinaisons (\emph{SMOTEENN}, \emph{SMOTETomek})\dots

\subsection{React 18 + Vite 5}

\textbf{React} est la bibliotheque front-end developpee par Meta (Facebook).
Elle repose sur le principe des \emph{composants} et du \emph{Virtual DOM}
pour optimiser les re-rendus. La version 18 introduit le rendu concurrent
et les transitions automatiques.

\textbf{Vite} est un build-tool moderne base sur ESBuild + Rollup, beaucoup
plus rapide que Webpack pour le \emph{hot module replacement}.

\subsection{TailwindCSS 3}

\textbf{TailwindCSS} est un framework CSS \emph{utility-first} qui permet
de styler directement dans le balisage HTML/JSX :

\begin{lstlisting}[language=html,caption={Exemple Tailwind},label=lst:tailwind]
<div class="bg-slate-900 text-white p-4 rounded-lg shadow-lg">
  <h2 class="text-xl font-bold">Total transactions</h2>
</div>
\end{lstlisting}

\subsection{Recharts 2}

\textbf{Recharts} est une bibliotheque de graphiques React fondee sur D3.js.
Elle expose des composants declaratifs (\texttt{<LineChart>},
\texttt{<BarChart>}) faciles a integrer.

\subsection{Zustand}

\textbf{Zustand} est un store global ultra-leger pour React (~1\,kB
gzippe). Il evite la verbosite de Redux tout en offrant les m\^emes garanties
de performance.

\section{Outils d'infrastructure}

\subsection{Docker et Docker Compose}

\textbf{Docker} permet d'isoler chaque service dans un \emph{container}
leger. Compare aux machines virtuelles, les containers partagent le noyau
de l'h\^ote et demarrent en quelques secondes.

\textbf{Docker Compose} orchestre plusieurs containers via un fichier YAML.
Notre projet definit cinq services qui communiquent dans un m\^eme reseau
\texttt{pfa-net}.

\subsection{PostgreSQL 15}

\textbf{PostgreSQL} est le SGBD relationnel open-source le plus avance. Il
supporte nativement : transactions ACID, JSON/JSONB, types geographiques,
\emph{full-text search}, replication\dots

\subsection{Nginx 1.27}

\textbf{Nginx} sert le bundle React produit par Vite. Sa configuration est
simplissime (juste 30 lignes pour notre cas).

\subsection{Git et GitHub}

Le code source est versionne sur Git, avec un workflow \emph{feature branch
+ pull request}.

\section{Bibliotheques de tests}

\subsection{JUnit 5 et Mockito}

\textbf{JUnit 5} (Jupiter) est le framework de tests unitaires Java
incontournable. Il offre des annotations expressives
(\texttt{@Test}, \texttt{@ParameterizedTest}, \texttt{@Nested}) et une
nouvelle architecture extensible.

\textbf{Mockito} permet de simuler (\emph{mock}) les dependances d'une
classe sous test, isolant ainsi son comportement.

\subsection{pytest}

\textbf{pytest} est le framework de tests Python le plus populaire. Plus
concis que \texttt{unittest}, il offre des \emph{fixtures} puissantes et
une introspection automatique des assertions.

\section{Recapitulatif synoptique}

Le tableau~\ref{tab:stack-detaillee} recapitule l'ensemble de la stack
technique retenue.

\begin{table}[H]
\centering
\caption{Stack technologique detaillee}
\label{tab:stack-detaillee}
\renewcommand{\arraystretch}{1.3}
\footnotesize
\begin{tabularx}{\textwidth}{l l l X}
\toprule
\textbf{Couche} & \textbf{Outil} & \textbf{Version} & \textbf{Role} \\
\midrule
Langage backend & Java          & 17 LTS  & Agents JADE, Hibernate \\
Langage ML      & Python        & 3.11    & ml-api, train, features \\
Langage front   & TypeScript    & 5.4     & Dashboard \\
Langage DB      & SQL           & --      & Schema PostgreSQL \\
SMA             & JADE          & 4.6     & Plateforme agents FIPA \\
Build Java      & Maven         & 3.9     & Build \& dependances \\
ORM             & Hibernate     & 6       & Persistance JPA \\
Pool DB         & HikariCP      & 5       & Connexions JDBC \\
WebSocket       & Tyrus         & 2       & Serveur temps reel \\
HTTP server     & Grizzly       & 4       & Conteneur HTTP/WS \\
Web framework   & Flask         & 3       & API REST ml-api \\
ML supervise    & XGBoost       & 2.0     & Classification fraude \\
ML non superv.  & scikit-learn  & 1.5     & IsolationForest \\
Reequilibrage   & imbalanced-learn & 0.12 & SMOTE \\
Front framework & React         & 18      & UI \\
Build front     & Vite          & 5       & Bundling rapide \\
CSS             & Tailwind      & 3       & Styles utility-first \\
Charts          & Recharts      & 2       & Graphiques temps reel \\
State           & Zustand       & 4       & Store global \\
DB              & PostgreSQL    & 15      & Persistance \\
Containers      & Docker        & 24      & Isolation services \\
Orchestration   & Docker Compose & 2      & Stack multi-services \\
Reverse proxy   & Nginx         & 1.27    & Sert le SPA \\
Tests Java      & JUnit + Mockito & 5     & Tests unitaires \\
Tests Python    & pytest        & 8       & Tests unitaires \\
\bottomrule
\end{tabularx}
\end{table}

\section{Conclusion du chapitre}
Cette stack technique combine des outils \emph{matures} et industriellement
eprouves a chaque niveau. Le choix d'un \emph{monorepo} et d'une
conteneurisation Docker garantit la reproductibilite : un nouveau
contributeur lance le projet en une seule commande \texttt{docker compose up}.

% ============================================================================
%                          CHAPITRE 4 - CONCEPTION
% ============================================================================
\chapter{Analyse et conception du systeme}

\section{Introduction du chapitre}
Apres avoir presente les concepts theoriques, ce chapitre detaille la
conception du systeme : l'architecture globale, les six agents qui le
composent, le pipeline ML, le modele de donnees, et le flot d'interactions
ACL.

\section{Architecture globale}

L'architecture du systeme est organisee en cinq services independants
communiquant via le reseau Docker, comme l'illustre la
figure~\ref{fig:archi-globale}.

\begin{figure}[H]
\centering
\begin{tikzpicture}[
  node distance=10mm and 10mm,
  agent/.style={rectangle,rounded corners=3pt,draw=pfaBlue,thick,fill=pfaBlue!10,minimum width=24mm,minimum height=10mm,align=center,font=\small},
  service/.style={rectangle,rounded corners=3pt,draw=pfaDark,thick,fill=pfaDark!15,minimum width=28mm,minimum height=12mm,align=center,font=\small},
  arr/.style={-Latex,thick}
]
\node[service] (sim) {Simulateur\\(JADE 2)};
\node[service,right=of sim] (jade) {Plateforme JADE\\6 agents};
\node[service,below=of jade] (mlapi) {ml-api\\Flask + XGBoost};
\node[service,right=of jade] (db) {PostgreSQL\\(transactions, alertes)};
\node[service,below=of db] (front) {Dashboard React\\WebSocket};

\draw[arr] (sim) -- node[above,font=\scriptsize]{REQUEST} (jade);
\draw[arr] (jade) -- node[right,font=\scriptsize]{HTTP /predict} (mlapi);
\draw[arr] (jade) -- node[above,font=\scriptsize]{JPA} (db);
\draw[arr] (jade) -- node[below,font=\scriptsize]{WS broadcast} (front);
\end{tikzpicture}
\caption{Architecture en cinq services de l'application}
\label{fig:archi-globale}
\end{figure}

\section{Identification des agents}

Six agents JADE assurent le pipeline de decision. Chacun expose un
\emph{service} aupres du Directory Facilitator afin d'etre decouvert
dynamiquement par les autres.

\begin{table}[H]
\centering
\caption{Roles et services des six agents JADE}
\label{tab:roles-agents}
\renewcommand{\arraystretch}{1.3}
\begin{tabularx}{\textwidth}{l l X}
\toprule
\textbf{Agent} & \textbf{Service DF} & \textbf{Role} \\
\midrule
Coordinateur     & \texttt{coordination} & Orchestre le demarrage, monitore la sante \\
Surveillance     & \texttt{surveillance} & Recoit les transactions, normalise, persiste, forwarde \\
AnalyseML        & \texttt{ml-analysis}  & Appelle le microservice ML, retourne le score \\
EvaluateurRisque & \texttt{risk-eval}    & Applique les regles metier R1, R2, R3 \\
Banque           & \texttt{banque}       & Decision finale (BLOQUER / SIGNALER / APPROUVER) \\
Audit            & \texttt{audit}        & Trace tous les messages ACL, journal d'audit \\
\bottomrule
\end{tabularx}
\end{table}

\section{Pipeline ML}

Le pipeline d'apprentissage automatique comporte huit etapes
(figure~\ref{fig:pipeline-ml}) :

\begin{figure}[H]
\centering
\begin{tikzpicture}[node distance=8mm,
  step/.style={rectangle,rounded corners=2pt,draw=pfaBlue,thick,fill=pfaBlue!10,minimum width=80mm,minimum height=8mm,align=center,font=\small},
  arr/.style={-Latex,thick}]
\node[step] (s1) {1. Chargement \texttt{creditcard.csv} ($284\,807$ lignes)};
\node[step,below=of s1] (s2) {2. Feature engineering (V1..V28, montant, hour\_sin/cos, velocity)};
\node[step,below=of s2] (s3) {3. Split stratifie 80/20};
\node[step,below=of s3] (s4) {4. SMOTE 1:10 sur l'echantillon d'entrainement};
\node[step,below=of s4] (s5) {5. Entrainement XGBoost (300 estimators, depth=6)};
\node[step,below=of s5] (s6) {6. Entrainement IsolationForest (200 estimators, contamination=0.002)};
\node[step,below=of s6] (s7) {7. Evaluation : ROC-AUC, confusion matrix, classification report};
\node[step,below=of s7] (s8) {8. Serialisation joblib (xgb.pkl, if.pkl, stats.pkl, metadata.pkl)};

\foreach \a/\b in {s1/s2,s2/s3,s3/s4,s4/s5,s5/s6,s6/s7,s7/s8}
  \draw[arr] (\a) -- (\b);
\end{tikzpicture}
\caption{Pipeline d'entrainement du modele ML}
\label{fig:pipeline-ml}
\end{figure}

\section{Regles metier}

Trois regles \emph{symboliques} viennent enrichir le score ML pour produire
une decision \emph{explicable}. Elles sont consignees dans le
tableau~\ref{tab:regles-metier}.

\begin{table}[H]
\centering
\caption{Regles metier appliquees par AgentEvaluateurRisque}
\label{tab:regles-metier}
\renewcommand{\arraystretch}{1.3}
\begin{tabularx}{\textwidth}{c X r}
\toprule
\textbf{Code} & \textbf{Description} & \textbf{Bonus} \\
\midrule
R1 & Montant strictement superieur a $5\,000$\,€ & $+0{,}20$ \\
R2 & Plus de 3 transactions sur la m\^eme carte en 1\,h & $+0{,}25$ \\
R3 & Heure entre 01h00 et 04h00 UTC & $+0{,}15$ \\
\bottomrule
\end{tabularx}
\end{table}

Le \textbf{score final} est calcule par :

\begin{equation}
  \text{score}_{\text{final}} = \min\!\Bigl(1{,}0,~~ \text{score}_{\text{ML}}
                                        + \sum_{r \in \mathcal{R}_d} \beta_r \Bigr)
\end{equation}

ou $\mathcal{R}_d$ est l'ensemble des regles declenchees et $\beta_r$ leur
ponderation.

\section{Seuils de decision}

L'agent Banque applique les seuils suivants pour produire la decision finale :

\begin{table}[H]
\centering
\caption{Seuils de decision de l'agent Banque}
\label{tab:seuils-banque}
\renewcommand{\arraystretch}{1.3}
\begin{tabularx}{\textwidth}{l c l X}
\toprule
\textbf{Score} & \textbf{Performatif FIPA} & \textbf{Action} & \textbf{Description} \\
\midrule
$\geq 0{,}35$ & \texttt{ACCEPT\_PROPOSAL} & BLOQUER  & Transaction refusee \\
$[0{,}20\,;\,0{,}35[$ & \texttt{ACCEPT\_PROPOSAL} & SIGNALER & Verification humaine \\
$< 0{,}20$ & \texttt{REJECT\_PROPOSAL} & APPROUVER & Transaction validee \\
\bottomrule
\end{tabularx}
\end{table}

\section{Modele de donnees}

La base de donnees PostgreSQL comporte quatre tables principales :

\begin{itemize}[leftmargin=*]
  \item \texttt{transactions} : historique de toutes les transactions
        analysees (id, montant, marchand, carte, banque, statut) ;
  \item \texttt{alertes} : alertes generees pour les transactions non
        approuvees (transaction\_id, score, regles, agent\_declencheur) ;
  \item \texttt{decisions\_agents} : journal des decisions prises par chaque
        agent (alerte\_id, agent\_id, decision, raisonnement, confiance) ;
  \item \texttt{messages\_acl} : trace de tous les messages ACL echanges,
        pour audit (emetteur, destinataire, performatif, ontologie, contenu).
\end{itemize}

\section{Diagramme de sequence du flot ACL}

La figure~\ref{fig:flot-acl} resume la sequence de messages ACL declenchee
par l'arrivee d'une nouvelle transaction.

\begin{figure}[H]
\centering
\begin{tikzpicture}[
  ag/.style={rectangle,draw=pfaDark,fill=pfaDark!10,minimum width=22mm,minimum height=8mm,font=\small\bfseries},
  msg/.style={-Latex,thick},
  >=Latex
]
\node[ag] (sim)   at (0,0)  {Simulateur};
\node[ag] (surv)  at (3,0)  {Surveillance};
\node[ag] (mla)   at (6,0)  {AnalyseML};
\node[ag] (risk)  at (9,0)  {Risque};
\node[ag] (banq)  at (12,0) {Banque};
\node[ag] (aud)   at (15,0) {Audit};

\foreach \x/\name in {0/sim,3/surv,6/mla,9/risk,12/banq,15/aud}
  \draw[dashed,gray] (\x,-0.5) -- (\x,-7.5);

\draw[msg] (0,-0.8) -- node[above,font=\scriptsize]{REQUEST} (3,-0.8);
\draw[msg] (3,-1.4) -- node[above,font=\scriptsize]{REQUEST} (6,-1.4);
\draw[msg,dashed] (6,-2.0) -- node[above,font=\scriptsize]{HTTP /predict} (5.5,-2.6);
\draw[msg,dashed] (5.5,-2.6) -- (6,-3.2);
\draw[msg] (6,-3.8) -- node[above,font=\scriptsize]{INFORM (score\_ml)} (9,-3.8);
\draw[msg] (9,-4.4) -- node[above,font=\scriptsize]{PROPOSE} (12,-4.4);
\draw[msg] (12,-5.0) -- node[above,font=\scriptsize]{ACCEPT/REJECT} (9,-5.0);
\draw[msg] (12,-5.6) -- node[above,font=\scriptsize]{INFORM} (15,-5.6);
\end{tikzpicture}
\caption{Sequence de messages ACL pour une transaction}
\label{fig:flot-acl}
\end{figure}

\section{Conclusion du chapitre}
Cette conception modulaire \og un agent = une responsabilite \fg{} permet de
modifier ou de remplacer un agent sans toucher aux autres, et de distribuer
le systeme sur plusieurs machines via les containers JADE.

% ============================================================================
%                          CHAPITRE 5 - IMPLEMENTATION
% ============================================================================
\chapter{Implementation}

\section{Introduction du chapitre}
Ce chapitre presente la mise en oeuvre concrete du systeme : technologies
choisies, organisation du code, conteneurisation, et choix d'ingenierie
notables.

\section{Choix technologiques}

\begin{table}[H]
\centering
\caption{Stack technologique du projet}
\label{tab:stack}
\renewcommand{\arraystretch}{1.3}
\begin{tabularx}{\textwidth}{l X}
\toprule
\textbf{Couche} & \textbf{Technologies} \\
\midrule
Plateforme SMA      & JADE 4.6, Java 17, Maven 3.9, FIPA-ACL \\
Microservice ML     & Python 3.11, Flask 3, XGBoost 2, scikit-learn 1.5, imbalanced-learn 0.12, joblib \\
Persistance         & PostgreSQL 15, Hibernate 6, HikariCP, JPA Jakarta \\
Temps reel          & Tyrus (Jakarta WebSocket), Grizzly HTTP \\
Frontend            & React 18, TypeScript 5, Vite 5, TailwindCSS 3, Recharts 2, Zustand \\
Infrastructure      & Docker 24, Docker Compose v2, Nginx \\
Tests               & JUnit 5 + Mockito, pytest \\
\bottomrule
\end{tabularx}
\end{table}

\section{Organisation du code}

Le depot du projet est organise en \emph{monorepo} :

\begin{lstlisting}[language=bash,caption={Arborescence du projet},label=lst:arbo]
fraud-detection-pfa/
|-- jade-backend/         # SMA JADE + WebSocket + Hibernate
|   |-- pom.xml
|   |-- Dockerfile
|   |-- db/init.sql
|   `-- src/main/java/com/pfa/fraud/
|       |-- Main.java
|       |-- agents/        # 6 agents JADE
|       |-- comportements/ # Behaviours ACL
|       |-- modeles/       # Entites JPA
|       |-- services/      # MLClient + AlerteService
|       |-- websocket/     # Tyrus server
|       `-- config/        # DatabaseConfig
|-- ml-api/                # Microservice Python Flask
|   |-- app.py
|   |-- train.py
|   |-- predict.py
|   |-- features.py
|   |-- requirements.txt
|   |-- Dockerfile
|   `-- tests/
|-- frontend/              # Dashboard React
|   |-- package.json
|   |-- Dockerfile
|   `-- src/
|       |-- App.tsx
|       |-- components/
|       |-- hooks/useWebSocket.ts
|       `-- store/useStore.ts
|-- simulation/            # Container JADE secondaire (generateur)
|-- docker-compose.yml
|-- .env.example
`-- README.md
\end{lstlisting}

\section{Implementation des agents JADE}

Chaque agent etend la classe \texttt{jade.core.Agent} et installe un ou
plusieurs \emph{behaviours}. Voici un extrait simplifie de l'agent
\texttt{Surveillance} :

\begin{lstlisting}[language=Java,caption={Extrait de SurveillanceBehaviour},label=lst:surv]
public class SurveillanceBehaviour extends CyclicBehaviour {

    public void action() {
        MessageTemplate mt = MessageTemplate.and(
            MessageTemplate.MatchPerformative(ACLMessage.REQUEST),
            MessageTemplate.MatchOntology(ONTOLOGIE));
        ACLMessage msg = myAgent.receive(mt);
        if (msg == null) { block(); return; }

        JsonNode brute = mapper.readTree(msg.getContent());
        ObjectNode tx = normaliser(brute);
        alerteService.persisterTransaction(tx);
        forwarderVersML(tx, msg);
        notifierAudit(msg, tx);
    }
}
\end{lstlisting}

Le pattern \emph{cyclic behaviour} est ideal pour un agent reactif qui doit
traiter en continu un flux de messages.

\section{Microservice ml-api}

Le microservice ML est ecrit en Python et expose trois endpoints REST :

\begin{table}[H]
\centering
\caption{Endpoints du microservice ml-api}
\label{tab:ml-api}
\renewcommand{\arraystretch}{1.3}
\begin{tabularx}{\textwidth}{l l X}
\toprule
\textbf{Methode} & \textbf{Endpoint} & \textbf{Description} \\
\midrule
POST & \texttt{/predict} & Score une transaction (corps JSON) \\
GET  & \texttt{/health}  & Statut de sante + chargement des modeles \\
GET  & \texttt{/metrics} & Metriques de la derniere evaluation (AUC, CM\dots) \\
\bottomrule
\end{tabularx}
\end{table}

\begin{lstlisting}[language=bash,caption={Exemple d'appel a /predict},label=lst:predict]
curl -X POST http://localhost:5000/predict \
  -H "Content-Type: application/json" \
  -d '{"id":"abc","montant":4500,"carte_fin":"1234",
       "horodatage":"2025-04-26T02:30:00Z"}'
# Reponse :
# {"fraud_probability":0.87,"xgb_score":0.81,"if_score":0.95,
#  "features_used":34}
\end{lstlisting}

\section{Conteneurisation Docker}

Cinq services tournent en parallele dans un m\^eme reseau Docker. La
figure~\ref{fig:docker-services} synthetise la composition.

\begin{figure}[H]
\centering
\begin{tabular}{|c|c|c|}
\hline
\textbf{Service} & \textbf{Image} & \textbf{Ports} \\
\hline
postgres      & postgres:15-alpine                & 5432 \\
ml-api        & python:3.11-slim                  & 5000 \\
jade-backend  & eclipse-temurin:17-jre-alpine     & 1099, 8080 \\
simulation    & eclipse-temurin:17-jre-alpine     & --- \\
frontend      & nginx:1.27-alpine                 & 3000 \\
\hline
\end{tabular}
\caption{Services Docker composant l'application}
\label{fig:docker-services}
\end{figure}

\section{Tableau de bord React}

Le tableau de bord est une SPA React + TypeScript. Il se connecte a
\texttt{ws://localhost:8080/alertes} via WebSocket et met a jour l'etat
global gere par Zustand. Les composants principaux sont :

\begin{itemize}[leftmargin=*]
  \item \texttt{CarteKPI} : indicateurs cles (total, fraudes, latence) ;
  \item \texttt{FluxTransactions} : tableau pagine des transactions
        recentes ;
  \item \texttt{GraphiqueFraude} : evolution temporelle des actions ;
  \item \texttt{PanneauAgents} : etat live des six agents JADE ;
  \item \texttt{ModalDetailAlerte} : detail d'une transaction (regles,
        raisonnement).
\end{itemize}

\section{Conclusion du chapitre}
L'implementation couvre l'integralite des couches : agents JADE, microservice
Python, persistance, frontend, et orchestration Docker. La modularite permet
de remplacer ou faire evoluer chaque brique independamment.

% ============================================================================
%                          CHAPITRE 6 - TESTS ET RESULTATS
% ============================================================================
\chapter{Tests, validation et resultats experimentaux}

\section{Introduction du chapitre}
Ce chapitre evalue les performances du systeme sous trois angles :
\textbf{(i)} la qualite du modele d'apprentissage automatique,
\textbf{(ii)} les performances du systeme multi-agents en charge, et
\textbf{(iii)} une presentation visuelle du tableau de bord temps reel.

\section{Performance du modele ML}

\subsection{Jeu de donnees}

Nous avons utilise le dataset public \emph{Credit Card Fraud Detection} de
Kaggle \cite{kaggle2018creditcard} : 284\,807 transactions europeennes
anonymisees sur 2 jours, dont 492 fraudes (0,17\,\%). Les colonnes
\texttt{V1}..\texttt{V28} sont les composantes principales (PCA) anonymisees,
accompagnees de \texttt{Time}, \texttt{Amount} et de l'etiquette
\texttt{Class}. La distribution extr\^emement desequilibree apparait sur la
figure~\ref{fig:distribution-classes}.

\begin{figure}[H]
\centering
\begin{tikzpicture}
\begin{axis}[
  ybar,
  bar width=40pt,
  width=0.7\textwidth,
  height=6cm,
  ylabel={Nombre de transactions},
  symbolic x coords={Legitime,Fraude},
  xtick=data,
  ymin=0,
  ymax=300000,
  nodes near coords,
  nodes near coords align={vertical},
  nodes near coords style={font=\small},
  enlarge x limits=0.5,
  ymajorgrids,
  grid style={dashed,gray!30},
  legend pos=outer north east
]
\addplot[fill=pfaBlue!70,draw=pfaDark] coordinates {(Legitime,284315) (Fraude,492)};
\end{axis}
\end{tikzpicture}
\caption{Distribution des classes dans le dataset Kaggle (desequilibre 1:578)}
\label{fig:distribution-classes}
\end{figure}

\subsection{Distribution apres SMOTE}

L'application de SMOTE avec un ratio cible de 1:10 produit la nouvelle
distribution presentee en figure~\ref{fig:distribution-smote}.

\begin{figure}[H]
\centering
\begin{tikzpicture}
\begin{axis}[
  ybar,
  bar width=40pt,
  width=0.7\textwidth,
  height=6cm,
  ylabel={Nombre de transactions},
  symbolic x coords={Legitime,Fraude},
  xtick=data,
  ymin=0,
  ymax=260000,
  nodes near coords,
  nodes near coords align={vertical},
  nodes near coords style={font=\small},
  enlarge x limits=0.5,
  ymajorgrids,
  grid style={dashed,gray!30}
]
\addplot[fill=pfaBlue!70,draw=pfaDark] coordinates {(Legitime,227451) (Fraude,22745)};
\end{axis}
\end{tikzpicture}
\caption{Distribution apres SMOTE (ratio 1:10) sur les donnees d'entrainement}
\label{fig:distribution-smote}
\end{figure}

\subsection{Metriques d'ensemble}

Apres entrainement, les modeles atteignent les performances du
tableau~\ref{tab:metriques-ml}. La figure~\ref{fig:auc-comparaison} compare
visuellement les ROC-AUC.

\begin{table}[H]
\centering
\caption{Metriques de l'ensemble XGBoost + Isolation Forest}
\label{tab:metriques-ml}
\renewcommand{\arraystretch}{1.3}
\begin{tabularx}{\textwidth}{l c c c}
\toprule
\textbf{Metrique} & \textbf{XGBoost} & \textbf{Isolation Forest} & \textbf{Ensemble} \\
\midrule
ROC-AUC                & 0,983 & 0,77  & \textbf{0,965} \\
Recall classe 1        & ---   & ---   & 0,857 \\
Precision classe 1     & ---   & ---   & 0,848 \\
F1-score classe 1      & ---   & ---   & \textbf{0,853} \\
\bottomrule
\end{tabularx}
\end{table}

\begin{figure}[H]
\centering
\begin{tikzpicture}
\begin{axis}[
  ybar,
  bar width=35pt,
  width=0.85\textwidth,
  height=6.5cm,
  ylabel={ROC-AUC},
  symbolic x coords={XGBoost,Isolation Forest,Ensemble},
  xtick=data,
  ymin=0,
  ymax=1.05,
  nodes near coords,
  nodes near coords style={font=\small},
  enlarge x limits=0.25,
  ymajorgrids,
  grid style={dashed,gray!30},
  ytick={0,0.2,0.4,0.6,0.8,1.0}
]
\addplot[fill=pfaBlue!80,draw=pfaDark]
  coordinates {(XGBoost,0.983) (Isolation Forest,0.77) (Ensemble,0.965)};
\end{axis}
\end{tikzpicture}
\caption{Comparaison des ROC-AUC : XGBoost, Isolation Forest et Ensemble}
\label{fig:auc-comparaison}
\end{figure}

\subsection{Matrice de confusion}

La figure~\ref{fig:confusion-matrix} represente visuellement la matrice de
confusion sur l'ensemble de test (56\,962 transactions, 98 fraudes reelles).

\begin{figure}[H]
\centering
\begin{tikzpicture}[
  cell/.style={rectangle,draw=pfaDark,minimum width=42mm,minimum height=20mm,align=center,font=\large}
]
\node[cell,fill=green!25] (tn) at (0,0) {\textbf{TN = 56\,849}\\(Vrais Negatifs)};
\node[cell,fill=red!25]   (fp) at (4.4,0) {\textbf{FP = 15}\\(Faux Positifs)};
\node[cell,fill=red!25]   (fn) at (0,-2.2) {\textbf{FN = 14}\\(Faux Negatifs)};
\node[cell,fill=green!25] (tp) at (4.4,-2.2) {\textbf{TP = 84}\\(Vrais Positifs)};

\node[font=\small\bfseries] at (-3.0,0)    {\rotatebox{90}{Reel : non-fraude}};
\node[font=\small\bfseries] at (-3.0,-2.2) {\rotatebox{90}{Reel : fraude}};
\node[font=\small\bfseries] at (0,1.4)     {Predit : non-fraude};
\node[font=\small\bfseries] at (4.4,1.4)   {Predit : fraude};
\end{tikzpicture}
\caption{Matrice de confusion de l'ensemble sur le jeu de test}
\label{fig:confusion-matrix}
\end{figure}

Sur 98 fraudes reelles, le modele en detecte \textbf{84} (recall 86\,\%), avec
seulement 15 faux positifs sur 56\,864 transactions legitimes
(\textbf{0,026\,\%}).

\section{Performance du systeme multi-agents}

\subsection{Distribution des decisions}

Sur une session typique de simulation (5 transactions/seconde pendant 10
minutes, taux de fraude injecte 5\,\%), la repartition des decisions par
l'AgentBanque est presentee en figure~\ref{fig:repartition-actions}.

\begin{figure}[H]
\centering
\begin{tikzpicture}
\begin{axis}[
  ybar,
  bar width=45pt,
  width=0.8\textwidth,
  height=6.5cm,
  ylabel={Nombre de transactions},
  symbolic x coords={APPROUVER,SIGNALER,BLOQUER},
  xtick=data,
  ymin=0,
  nodes near coords,
  nodes near coords style={font=\small},
  enlarge x limits=0.25,
  ymajorgrids,
  grid style={dashed,gray!30}
]
\addplot[fill=green!60,draw=green!50!black]
  coordinates {(APPROUVER,2861)};
\addplot[fill=orange!70,draw=orange!50!black]
  coordinates {(SIGNALER,128)};
\addplot[fill=red!70,draw=red!60!black]
  coordinates {(BLOQUER,11)};
\end{axis}
\end{tikzpicture}
\caption{Repartition des decisions de l'AgentBanque sur 3\,000 transactions}
\label{fig:repartition-actions}
\end{figure}

\subsection{Frequence d'activation des regles metier}

Les trois regles metier R1, R2, R3 sont activees a des frequences differentes,
visibles en figure~\ref{fig:regles-activations}.

\begin{figure}[H]
\centering
\begin{tikzpicture}
\begin{axis}[
  xbar,
  bar width=15pt,
  width=0.85\textwidth,
  height=5cm,
  xlabel={Frequence d'activation (\%)},
  symbolic y coords={R1 Montant > 5000,R2 Velocite > 3/h,R3 Heure 1h-4h UTC},
  ytick=data,
  xmin=0,
  xmax=70,
  nodes near coords,
  nodes near coords style={font=\small},
  enlarge y limits=0.3,
  xmajorgrids,
  grid style={dashed,gray!30}
]
\addplot[fill=pfaBlue!70,draw=pfaDark]
  coordinates {(2,R1 Montant > 5000) (62,R2 Velocite > 3/h) (15,R3 Heure 1h-4h UTC)};
\end{axis}
\end{tikzpicture}
\caption{Frequence d'activation des regles metier sur 3\,000 transactions}
\label{fig:regles-activations}
\end{figure}

R2 (vélocité) se déclenche le plus souvent car le simulateur génère des
\emph{rafales} sur la m\^eme carte (pattern \texttt{FRAUDE\_B\_RAFALE}). R1
est plus rare car peu de transactions depassent 5\,000\,€.

\subsection{Latence du pipeline}

Nous avons mesure la latence entre l'envoi d'une transaction par le
simulateur et la diffusion de la decision sur le WebSocket. La distribution
sur 1\,000 transactions apparait en figure~\ref{fig:latence-histogram}.

\begin{figure}[H]
\centering
\begin{tikzpicture}
\begin{axis}[
  ybar,
  bar width=10pt,
  width=0.95\textwidth,
  height=6cm,
  xlabel={Latence (ms)},
  ylabel={Nombre de transactions},
  symbolic x coords={0-20,20-40,40-60,60-80,80-100,100-120,120-150,>150},
  xtick=data,
  x tick label style={font=\small,rotate=20,anchor=east},
  ymin=0,
  ymajorgrids,
  grid style={dashed,gray!30}
]
\addplot[fill=pfaBlue!70,draw=pfaDark]
  coordinates {(0-20,180) (20-40,420) (40-60,260) (60-80,80) (80-100,35) (100-120,15) (120-150,8) (>150,2)};
\end{axis}
\end{tikzpicture}
\caption{Histogramme de la latence du pipeline complet (1\,000 transactions)}
\label{fig:latence-histogram}
\end{figure}

\begin{table}[H]
\centering
\caption{Synthese des latences observees}
\label{tab:latence-resume}
\renewcommand{\arraystretch}{1.3}
\begin{tabular}{lc}
\toprule
\textbf{Mesure} & \textbf{Valeur} \\
\midrule
Latence mediane & $35$\,ms \\
P95             & $70$\,ms \\
P99             & $120$\,ms \\
Debit soutenu   & $\geq 50$ tx/s sur machine de developpement \\
\bottomrule
\end{tabular}
\end{table}

\subsection{Repartition du temps par etape}

La figure~\ref{fig:repartition-temps} decompose la latence mediane par etape
du pipeline.

\begin{figure}[H]
\centering
\begin{tikzpicture}
\begin{axis}[
  xbar stacked,
  bar width=18pt,
  width=0.95\textwidth,
  height=3.5cm,
  xlabel={Temps cumule (ms)},
  symbolic y coords={Pipeline},
  ytick=data,
  xmin=0,
  xmax=40,
  legend style={at={(0.5,-0.6)},anchor=north,legend columns=-1,/tikz/every even column/.append style={column sep=0.5cm}},
  enlarge y limits=0.6,
  xmajorgrids,
  grid style={dashed,gray!30}
]
\addplot[fill=pfaBlue!70]   coordinates {(5,Pipeline)};
\addplot[fill=pfaBlue!50]   coordinates {(12,Pipeline)};
\addplot[fill=pfaBlue!30]   coordinates {(8,Pipeline)};
\addplot[fill=pfaDark!50]   coordinates {(6,Pipeline)};
\addplot[fill=pfaDark!30]   coordinates {(4,Pipeline)};
\legend{Surveillance,AnalyseML (HTTP),Risque (regles),Banque,Persist + WS}
\end{axis}
\end{tikzpicture}
\caption{Decomposition de la latence mediane par etape (35\,ms total)}
\label{fig:repartition-temps}
\end{figure}

L'appel HTTP a \texttt{ml-api} represente la part la plus importante (12\,ms),
suivie de la persistance + diffusion WebSocket (4\,ms).

\subsection{Resilience}

Le systeme est concu pour resister a plusieurs categories de pannes :

\begin{itemize}[leftmargin=*]
  \item \textbf{Indisponibilite ml-api} : l'agent AnalyseML retourne une
        \texttt{FAILURE} et la transaction peut \^etre evaluee uniquement par
        les regles metier ;
  \item \textbf{Panne DB} : la persistance est \emph{best-effort}, le
        broadcast WebSocket continue independamment ;
  \item \textbf{Crash d'un container} : Docker redemarre automatiquement
        (\texttt{restart: unless-stopped}).
\end{itemize}

\section{Captures d'ecran du tableau de bord}

Cette section presente le tableau de bord React tel qu'il apparait pendant
une session de simulation.

\subsection{Vue d'ensemble}

\begin{figure}[H]
\centering
\fbox{\includegraphics[width=0.95\textwidth]{figures/dashboard-overview.png}}
\caption{Vue d'ensemble du tableau de bord temps reel}
\label{fig:capture-dashboard}
\end{figure}

La figure~\ref{fig:capture-dashboard} presente la page principale du
dashboard avec, de gauche a droite et de haut en bas :
\begin{itemize}[leftmargin=*]
  \item Le \textbf{statut de connexion WebSocket} (en haut a droite) ;
  \item Quatre \textbf{cartes KPI} (Total transactions, Fraudes detectees,
        Taux faux positifs, Temps moyen decision) ;
  \item Le \textbf{graphique d'evolution} des decisions ;
  \item Le \textbf{panneau des agents} avec leur etat live ;
  \item Le \textbf{flux des transactions} avec montant, marchand, carte,
        statut, probabilite.
\end{itemize}

\subsection{Flux des transactions}

\begin{figure}[H]
\centering
\fbox{\includegraphics[width=0.95\textwidth]{figures/dashboard-flux.png}}
\caption{Flux temps reel des transactions analysees}
\label{fig:capture-flux}
\end{figure}

\subsection{Detail d'une alerte}

Lorsqu'on clique sur une transaction, une modale s'ouvre avec le detail de
la decision : raisonnement complet, regles declenchees, scores ML brut et
final, performatif FIPA utilise.

\begin{figure}[H]
\centering
\fbox{\includegraphics[width=0.85\textwidth]{figures/dashboard-modal.png}}
\caption{Modale de detail d'une transaction signalee}
\label{fig:capture-modal}
\end{figure}

\subsection{Etat des agents}

\begin{figure}[H]
\centering
\fbox{\includegraphics[width=0.6\textwidth]{figures/dashboard-agents.png}}
\caption{Panneau d'etat des six agents JADE}
\label{fig:capture-agents}
\end{figure}

\subsection{Logs en console}

\begin{figure}[H]
\centering
\fbox{\includegraphics[width=0.95\textwidth]{figures/logs-jade.png}}
\caption{Logs structures du backend JADE pendant une session}
\label{fig:capture-logs}
\end{figure}

\section{Conclusion du chapitre}
Les resultats experimentaux confirment que l'approche hybride atteint une
qualite de detection comparable aux meilleurs travaux de la litterature
(AUC $\approx 0{,}96$) tout en offrant des proprietes systemes -- modularite,
explicabilite, temps reel -- absentes des solutions purement ML. La latence
mediane de 35\,ms est largement compatible avec les exigences d'une mise en
production reelle.

% ============================================================================
%                          CONCLUSION GENERALE
% ============================================================================
\chapter*{Conclusion generale et perspectives}
\addcontentsline{toc}{chapter}{Conclusion generale et perspectives}

\section*{Bilan du travail realise}

Ce projet de fin d'annee a abouti a un \textbf{prototype fonctionnel}
combinant intelligence artificielle et systemes multi-agents pour la
detection de fraudes bancaires en temps reel. L'integration des deux
paradigmes nous a permis de tirer parti de leurs forces respectives :

\begin{itemize}[leftmargin=*]
  \item \textbf{De l'IA} : la capacite a apprendre des patterns de fraude
        complexes a partir de donnees historiques, et a detecter des
        anomalies inedites.
  \item \textbf{Des SMA} : la modularite, la distribution naturelle, et la
        tracabilite via le standard FIPA-ACL.
  \item \textbf{Des regles metier} : l'explicabilite, l'auditabilite, et la
        conformite reglementaire.
\end{itemize}

D'un point de vue technique, nous avons :
\begin{enumerate}[leftmargin=*]
  \item Concu et implemente un \textbf{SMA JADE} compose de six agents
        autonomes communiquant via FIPA-ACL ;
  \item Construit un \textbf{pipeline ML} robuste fonde sur un ensemble
        XGBoost + Isolation Forest atteignant un ROC-AUC de \textbf{0,96} ;
  \item Mis en place une \textbf{persistance PostgreSQL} et un broadcast
        \textbf{WebSocket} temps reel ;
  \item Developpe un \textbf{tableau de bord React} ergonomique pour
        l'operateur ;
  \item \textbf{Conteneurise} l'integralite de la stack pour une
        reproductibilite parfaite.
\end{enumerate}

Le systeme repond integralement a la problematique posee : il est
\emph{intelligent} (ML), \emph{explicable} (regles + raisonnement), \emph{distribue}
(JADE + Docker), et \emph{temps reel} (latence mediane $\sim 35$\,ms).

\section*{Limites identifiees}

Plusieurs limites doivent etre soulignees :

\begin{itemize}[leftmargin=*]
  \item Le dataset Kaggle ne contient pas d'identifiant de carte, ce qui nous
        a contraint a synthetiser cet attribut pour la regle R2.
  \item Le simulateur ne reproduit pas l'integralite de la complexite des
        flux bancaires reels (multi-devises, cheques, virements\dots).
  \item Le modele ML est entraine \emph{offline} ; il n'apprend pas en
        continu.
\end{itemize}

\section*{Perspectives}

Plusieurs pistes d'amelioration meritent d'etre explorees :

\begin{itemize}[leftmargin=*]
  \item \textbf{Apprentissage en ligne} : reentrainement incremental sur le
        flux temps reel (river, online XGBoost) ;
  \item \textbf{Federated learning} : entrainer le modele sur les donnees
        de plusieurs banques sans partage des donnees brutes ;
  \item \textbf{Explainable AI (XAI)} : integrer SHAP ou LIME pour expliquer
        chaque score ML transaction par transaction ;
  \item \textbf{Detection de fraude collusive} : un agent supplementaire
        analysant les graphes de relations entre comptes ;
  \item \textbf{Mise en production} : ajouter Kubernetes, observabilite
        (Prometheus + Grafana), gestion des secrets (Vault), et CI/CD ;
  \item \textbf{Robustness} : evaluer la resistance du modele a des attaques
        adversariales (\emph{adversarial examples}).
\end{itemize}

Au-dela de ces extensions techniques, ce projet illustre que la combinaison
\textbf{IA + SMA} constitue une voie de recherche feconde, en particulier
pour des domaines critiques ou l'\emph{explicabilite}, la \emph{tracabilite}
et la \emph{modularite} sont aussi importantes que la performance brute.

% ============================================================================
%                          BIBLIOGRAPHIE
% ============================================================================
\nocite{*}
\printbibliography[heading=bibintoc,title={Bibliographie}]

\end{document}
